\section{A Training Recipe for NaCl}

The application of NaCl to the \textsc{Lemaître} dataset presents a 
different challenge from the training on SNfactory data. The \textsc{Lemaître} 
sample contains a large number of supernovae and a correspondingly large 
number of model parameters. In this setting, fitting all 
parameters simultaneously from arbitrary initial values can lead to an 
inefficient training procedure. 

In particular, we found that when the initial model parameters are far from 
their optimal values, the error model can absorb part of the discrepancies 
between the model and the data. Although this provides a valid direction for 
the minimizer and can lead to a decrease in $Q(\boldsymbol{\beta})$, it can prevent the 
physical model parameters from being properly constrained during the initial 
stages of the fit. The non-linear nature of the model can also lead the 
minimizer towards undesirable local minima depending on the initial 
parameters. \\

To address these issues, I developed a sequential training procedure that 
provides increasingly informed initial estimates for the different components 
of the model before performing the full fit. This approach reduces the 
sensitivity of the minimization to the initial conditions, improves the 
stability of the training, and reduces the computational cost of the full 
optimization. The resulting training recipe is divided into five steps: 

\begin{enumerate}
    \item \underline{Light curve fit with the initial NaCl model:} 
    Only $X_0$, $X_1$, $c$, and $t_{max}$ are fitted for each 
    supernova, while the global model parameters are kept fixed.

    \item \underline{$SpecRecal$ fit with the initial NaCl model:} 
    The coefficients describing the spectroscopic recalibration are 
    estimated for each spectrum. Since the recalibration model is linear 
    in these coefficients, this step is computationally inexpensive. During this step, all other parameters are fixed at the best value found in the previous step. 

    \item \underline{Initial model guess:} 
    The parameters describing the model surfaces $M_0$, $M_1$, and $CL$ are 
    fitted alongside $X_0$, $X_1$, $c$, $t_{max}$ and $SpecRecal$ with a fixed error model. 
    This is done to obtain an initial model that is sufficiently close to the 
    solution and satisfies the required model constraints.

    \item \underline{Initial error model guess:} 
    With the spectrophotometric model fixed, the parameters of the error 
    model are fitted while the all of the other parameters are fixed to the value found in the previous step. 
    This is done to provide an initial estimate of the error model.

    \item \underline{Full training:} 
    All model, error model, and supernova parameters are released and 
    simultaneously fitted.
\end{enumerate}

The first two steps are computationally inexpensive. The initial model fit 
in step 3 is more demanding because of the large number of parameters 
associated with the $M_0$ and $M_1$ surfaces. Since the purpose of this step 
is only to obtain a suitable starting point for the subsequent optimization, 
the number of minimization iterations is limited to 15. The error-model fit 
in step 4 is again relatively inexpensive. The final step, in which all 
parameters are simultaneously released, constitutes the most computationally 
expensive part of the training.

%\subsection{Regularization}

%To try and overcome this issue, I've developed during this PhD a regularization scheme that adapts the force of the regularization based on the number of 
%points seen by each model parameter. This regularization scheme is inspired by \cite{lenz2022adaptiveregularizationbsplinemodels}. We can separate $\mathbf{P}$ matrix as 
%follows: 

%\begin{equation}
%   \mathbf{P} = \mathbf{P}^{data} + \mathbf{P}^{smooth}
%\end{equation}

%with $\mathbf{P}^{data}$ and $\mathbf{P}^{smooth}$ the matrices determining the strength based on the number of points and to smoothen the model where there is no data as 
%much as possible. In order to determine $\mathbf{P}^{data}$, we define $\mathcal{N}_i$ as the number of data points seen by each $\theta_{i}$ parameters ($\theta_{kl}$ written 
%as a vector). We then introduce for each $\mathcal{N}_i$ a parameter $\mu_i$ which will act as a gate function of $\mathcal{N}_i$. This will determine the force of our 
%regularization. The idea is that if $\mathcal{N}_i=0$ we regularize model strongly, otherwise we don't.\\

%$\mu_i$ is then defined as: 

%\begin{equation}
% \mu_i^{data} = \mathrm{max}(\mu_{reg}^{data} - \mathcal{N}_i, 0)
% \label{eqn:reg_mu}
%\end{equation}

%We can use that to define a diagonal matrix $\mathbf{M}^{data}$ where $\mathbf{M}_{ii}^{data} = \mu_i^{data}$. It is then clear that the simplest way to define 
%$\mathbf{P}^{data}$ will be to write it as:

%\begin{equation}
%   \mathbf{P}^{data} = \mathbf{M}^{data} \times \mathbb{1}
%\end{equation}

%\noindent We then define $\mathbf{P}^{smooth}$ as a penalty on the first order derivative of the model in phase, like in equation \ref{eqn:old_reg}. We apply the same 
%adaptive method mentioned previously and we get:

%\begin{equation}
%   \mathbf{P}^{smooth} = \mathbf{M}^{smooth} \times \mathbb{P}
%\end{equation}

%with $\mathbb{P}$ being : 
%\begin{equation}
%    \mathbb{P} = \begin{bmatrix} 
%    2  & -1& 0  &   &     &...  \\
%    -1 & 1 & -1 & 0 & ... & \\  
%    0  &   &    &...&     &  \\  
%       &   &... &   &     &  0\\ 
%       &...& 0  & -1& 1   & -1 \\  
%    ...&   &    & 0 & -1  & 2 \end{bmatrix}
%\end{equation}
%with $\mu_{reg}^{smooth}$ would be the hyperparameter determining the strength of the regularization on the derivatives.\\
%Since the hyperparameters are now included in the  $\mathbf{P}$ matrix, the penalty now becomes : 
%\begin{equation}
%    \boxed{\chi^2_{reg} = \beta^T \mathbf{P} \beta}
%    \label{eqn:chi2_reg_new}
%\end{equation}

%This adaptive regularization however presented issues when applied to supernova data. The minimization algorithm would not fully converge when this regularization is used. Moreover, 
%the recovered supernova parameters presented biases. 
